% GDZ PPN236018647,native532–534upper;HKUST533–535.
\clearpage\SourceLeaf{532}\ArithmeticHead{503}{R}{1796 SEPT. 6—1796 SEPT. 17.}
\NumberedItem{34}
\indent Methodus facilis inveniendi aeq[uationem] in \(y\) ex aeq[uatione] in \(x\), si ponatur:
\[
x^n+ax^{n-1}+bx^{n-2}\ \ldotp\,\ldotp\,\ldotp=y.
\]
\noindent\hfill[1796 Sept. 16.]

\medskip\begingroup\small
\indent \textsc{Gauss} hat wohl die sogenannte \textsc{Tschirnhausen}sche Transformation im Auge: Aus einer Gleichung: \(x^m+Ax^{m-1}+Bx^{m-2}+\ldotp\,\ldotp\,\ldotp+M=0\) soll die transformierte Gleichung in \(y\) abgeleitet werden, wenn \(y=x^n+ax^{n-1}+bx^{n-2}+\ldotp\,\ldotp\,\ldotp\) ist. Es könnte dann mit der \textls[120]{bequemen Methode}, die Gleichung in \(y\) zu finden, die mittels der Potenzsummen gemeint sein, die für den besondern Fall \(y=x^\tau\) im art. 348. der \textit{Analysis Residuorum}, Werke II, S. 223 als \textls[120]{Solutio prima} auseinandergesetzt ist. Man hätte so eine unmittelbare Anknüpfung an die Notiz Nr. 28 vom 21. August 1796.
\par\noindent\hfill\textsc{Loewy.}
\par\endgroup

\NumberedItem{35}
\indent Fractiones quarum denominator continet quantitates irrationales (quomodocunque?) in alias transmutare ab hoc incommodo liberatas.
\par\noindent\hfill[1796] Sept. 16.

\medskip\begingroup\small
\indent Vermutlich handelt es sich um die Umwandlung einer gebrochenen rationalen Funktion einer Wurzel einer algebraischen Gleichung in eine ganze rationale Funktion dieser Wurzel. Ausführlich behandelt \textsc{Gauss} die »principia talis transformationis, quum in libris algebraicis non inveniantur« im art. 11. der \textit{Methodus nova integralium valores per approximationem inveniendi} (Werke III, S. 177). Auf die Möglichkeit einer solchen Umwandlung wird auch schon im art. 360. III der \textit{Disquisitiones arithmeticae} (Werke I, S. 453) hingewiesen.
\par\noindent\hfill\textsc{Loewy.}
\par\endgroup

\NumberedItem{36}
\indent Coefficientes aeq[uationis] auxiliariae eliminationi inservientis ex radicibus aeq[uationis] datae determinati.
\par\noindent\hfill eod[em die, 1796 Sept. 16.]

\medskip\begingroup\small
\indent Vermutlich hängt diese Aufzeichnung mit den Untersuchungen über Elimination zusammen, von denen in der Aufzeichnung Nr. 89 des \textit{Tagebuchs} vom Juni 1798 die Rede ist.
\par\noindent\hfill\textsc{Loewy.}
\par\endgroup

\NumberedItem{37}
\indent Nova methodus qua resolutionem aequationum universalem investigare forsitanque invenire licebit. Scil[icet] transm[utetur] aeq[uatio] in aliam, cuius
\clearpage\SourceLeaf{533}\ArithmeticHead{504}{L}{TAGEBUCH MIT ERLÄUTERUNGEN.}
\noindent radices
\[
\alpha\rho'+\beta\rho''+\gamma\rho'''+\ldotp\,\ldotp\,\ldotp,
\]
\noindent ubi
\[
{}^{n}\!\surd1=\alpha,\ \beta,\ \gamma\ \text{etc.}
\]
\noindent et \(n\) numerus aequationis gradum denotans.
\par\noindent\hfill[1796] Sept. 17.

\medskip\begingroup\small
\indent Diese Aufzeichnung zeigt, daß \textsc{Gauss} mit Hilfe der sogenannten \textsc{Lagrange}schen Resolvente\footnote{Siehe \textsc{Lagrange}, \textit{Réflexions sur la résolution algébrique des équations}, Nouveaux Mémoires de l’Académie, année 1770, Berlin 1772, S. 134 und \textit{Suite des réflexions etc.}, ebenda, année 1771, Berlin 1773, S. 138, besonders S. 163 der \textit{Suite}, Oeuvres III, S. 331. Eine deutsche Übersetzung beider Teile dieser Abhandlung findet sich in \textit{Leonh. Eulers Einleitung in die Analysis des Unendlichen}, . . . übersetzt und mit . . . Zusätzen begleitet von J. A. Chr. \textsc{Michelsen}, 3. Buch: \textit{Die Theorie der Gleichungen aus den Schriften der Hrn. Euler und de la Grange}, Berlin 1791, S. 271 ff.; der zweite Teil beginnt auf S. 378. \textsc{Gauss} hat das \textsc{Michelsen}sche \textit{dritte Buch} spätestens im April 1797 gekannt, da er es in der Abhandlung \textit{Neuer Beweis des Lagrangeschen Lehrsatzes} (Werke VIII, S. 76) anführt, die er etwa im April 1797 (nämlich dritthalb Jahre vor dem 8. Oktober 1799, siehe den Brief an \textsc{Hindenburg} von diesem Tage, oben S. 429) durch \textsc{Kaestner} hat an \textsc{Hindenburg} senden lassen; vergl. auch die Tagebuchnotiz Nr. 49 vom 27. Dezember 1796. Die Arbeit von \textsc{Lagrange}, \textit{Réflexions sur la résolution etc.}, von der hier die Rede ist, meint \textsc{Pfaff} in dem oben S. 101 abgedruckten Briefe, wenn er S. 104, Zeile 5, 4 v. u. von »der früheren Abhandlung von \textsc{La Grange}, die bekanntlich auch von \textsc{Michelsen} deutsch übersetzt ist« spricht.}\ die algebraische Auflösbarkeit der allgemeinen algebraischen Gleichung untersucht und damals noch an die Möglichkeit gedacht hat, die Auflösung auf diese Weise auch wirklich zu finden. Daß \textsc{Gauss} sich viel mit dieser Frage beschäftigt hat, geht aus einer Stelle des nicht veröffentlichten Teiles des \textls[120]{Caput sextum} der \textit{Analysis Residuorum}\begingroup\renewcommand{\thefootnote}{**)}\footnote{Über diese Handschrift (Ea 9, Kapsel 40) vergleiche man die Bemerkungen, die R. \textsc{Dedekind} dem Abdruck zweier Abschnitte, Werke II, S. 240 hinzugefügt hat, und den Artikel 2 von P. \textsc{Bachmanns} Aufsatz »Über \textsc{Gauss}' zahlentheoretische Arbeiten«, Werke X 2, S. 6; danach ist das 6. Kapitel der \textit{Analysis Residuorum} vermutlich vor dem 20. Juli 1797 verfaßt worden.}\endgroup\ hervor, wo es im art. 262. heißt:
\par\endgroup

\indent Post illustrissimorum Geometrarum labores repetitos nulla spes superesse videtur Aequationum solutionem generalem possibilem esse (i. e. reductionem ad aequationes puras\textsuperscript{\(\dagger\))}). Sed tamen maxime est memorabile, aequationes omnes, ad quas solutio aequationis \(x^n=1\) ducit, resolvi sive ad puras eiusdem gradus reduci posse . . . .
\par\smallskip\noindent\rule{.24\linewidth}{.4pt}
\par\indent \(\dagger\)) Nos etiam huic rei multum operae impendimus; et tantum non de impossibilitate sumus certi. Forsan, quae in hoc genere meditati sumus et quae forsan ad demonstrationem rigorosam huius impossibilitatis ducere possunt, alia occasione publici iuris faciemus.

\clearpage\SourceLeaf{534}\ArithmeticHead{505}{R}{1796 SEPT. 17.—1796 OCTOB. 1.}
\begingroup\small
\indent Diese Stelle zeigt auch, daß \textsc{Gauss} schon im Jahre 1797 zu der Überzeugung gelangt war, die allgemeine algebraische Gleichung sei durch Wurzelziehen nicht auflösbar; öffentlich ausgesprochen hat er diese Überzeugung erst im art. 9. der \textit{Inauguraldissertation} (1799, Werke III, S. 17) und im art. 359. der \textit{Disquisitiones arithmeticae} (1801, Werke I, S. 449).

\indent Die Anwendung der \textsc{Lagrange}schen Resolvente auf die Kreisteilungsgleichung (siehe \textit{Disquisitiones arithmeticae} art. 360., Werke I, S. 450) ist erst in den Tagebuchnotizen Nr. 65 und Nr. 66 vom Juli 1797 angezeigt.
\par\noindent\hfill\textsc{Loewy.}
\par\endgroup
